% EM structure and teaching revision 20260928
\subsection{充电过程为何不能只算末态储能}
电容通过电阻 $R$ 接到恒压源 $U_0$，初始电压为零。由 $U_0=Ri+u_C$ 与 $i=C\dd u_C/\dd t$ 得
\[
 u_C=U_0(1-e^{-t/(RC)}),\qquad i=\frac{U_0}{R}e^{-t/(RC)}.
\]
电源总供能为 $\int_0^\infty U_0i\dd t=CU_0^2$，电容末态只储存 $CU_0^2/2$，另半转化为焦耳热。减小 $R$ 会增大峰值电流、缩短过程，却不改变理想模型中总焦耳热。
两个初始电压为 $U_1,U_2$ 的电容 $C_1,C_2$ 同名端相连，且体系无外接电源时，总电荷守恒给出
\[
 U_f=\frac{C_1U_1+C_2U_2}{C_1+C_2},\qquad
 W_i-W_f=\frac12\frac{C_1C_2}{C_1+C_2}(U_1-U_2)^2.
\]
能量减少量不会因为连接线画成理想导线就自动消失。实际过程还可能涉及导线电感、电磁辐射及介质损耗；若严格设为无损LC系统，则会振荡而不自行停在上述平衡末态。先说明终态能否建立，再使用末态电荷约束。

